\chapter*{Abstract}
\thispagestyle{plain}
\addcontentsline{toc}{chapter}{\numberline{}Abstract}
Software-Defined Networking (SDN) moves the control logic of a network into a logically centralised controller. This gives operators a global view of the traffic, but it also makes the controller and the switches attractive targets for Distributed Denial of Service (DDoS), Denial of Service (DoS), port scanning and brute-force attacks. Conventional machine-learning intrusion detection systems score every network flow independently as a row of a table. They therefore ignore \emph{who talks to whom}: a spoofed DDoS appears as thousands of harmless-looking one-packet flows, and a port scan appears as one host contacting many ports. Most published systems are also evaluated only offline on features such as CICFlowMeter statistics that a real OpenFlow switch cannot report, so they cannot be deployed in the network they were designed to protect.\\

The objective of this project is to design, implement and evaluate an intrusion detection and automated mitigation system for SDN that (i) represents the traffic of each time window as a host-to-host graph, (ii) uses a Graph Neural Network (GNN) to decide \emph{whether} the window contains an attack, \emph{which type} of attack it is, and \emph{which hosts} are the attackers, and (iii) closes the loop by having the SDN controller block or rate-limit the attackers automatically. The system is restricted to features that an OpenFlow~1.3 switch actually reports, so the same model runs on the datasets and on live flow statistics.\\

Traffic from the InSDN dataset (343,888 cleaned flows) is converted into 17,780 graph windows in which hosts are nodes and flows are edges, with 16 edge features and 16 node features derived from OpenFlow counters. A multi-task Graph Attention Network (GAT) built with PyTorch Geometric produces a graph-level attack-type prediction and a node-level attacker score. It is compared against Random Forest and XGBoost baselines under an identical time-ordered split, and against GCN, GraphSAGE and window-size ablations; PCA, SVD, ICA and filter, wrapper and embedded feature-selection methods are also compared. The trained model is exported with TorchScript and served by a FastAPI service. An os-ken controller polls flow statistics every two seconds from Open vSwitch in a Mininet/Docker laboratory and installs the drop or meter-based rate-limit rules returned by the service. On the held-out test split the GAT reaches a window-level F1-score of 0.998 with no false positives and an attacker-host F1-score of 0.998. The complete service path takes a median of 6.5~ms per window on a CPU. In the live laboratory, 12 of 12 attack runs were mitigated with a mean of 94.7\% of the attack traffic dropped and no benign host blocked.\\

The main contributions are an OpenFlow-compatible feature schema shared by the datasets and the live controller, a multi-task graph attention network that names both the attack type and the attacking hosts, a leakage-free evaluation protocol (time-ordered split with a gap, direction-corrected labels, benign overlay), a systematic comparison of graph layers, window sizes and eight feature-selection and reduction methods, and an end-to-end mitigation loop with measured outcomes. The evaluation also reports where the approach is limited: the rare attack classes are not learned, the model needs fine-tuning on laboratory traffic before it can be used live, legitimate clients are heavily throttled during a spoofed DDoS, and the time to mitigation (median 4.5~s) is bounded by the polling interval and not by the model.\\

\textbf{Keywords:} Software-Defined Networking, Intrusion Detection, Graph Neural Network, Graph Attention Network, OpenFlow, Mininet, Automated Mitigation, InSDN.
